\documentclass[12pt]{article}
\usepackage[utf8]{inputenc}
\usepackage[T1]{fontenc}
\usepackage{hyperref}
\title{Sardinian, Sassarese Language Resources for K-12 Curriculum Development}
\author{Compiled by Claude Sonnet 5.5 (claude-sonnet-5-5) for Andr\'as Kornai}
\date{}
\begin{document}
\maketitle

\section{Introduction}
Sassarese (ISO 639-3 sdc; Glottocode sass1235) is a Romance variety of north-western Sardinia centred on Sassari, about 100,000 speakers, rooted in medieval Tuscan with strong Sardinian influence and close to Gallurese and Corsican. It has no official regulating body. The resources found are modest: a learner course, dictionary samples and reference pages.

\section{Primary Educational Resources}
\begin{itemize}
\item A 50-lesson PDF course by Antoninu Rubattu, intended for new learners and elementary school students, with Italian used for examples and translations. It is catalogued by the Endangered Languages Project (availability of the file not verified): \url{https://endangeredlanguages.com/node/108376}
\end{itemize}
No official school curriculum or textbook was found.

\section{Linguistic and Reference Materials}
\begin{itemize}
\item Wikipedia article on Sassarese: \url{https://en.wikipedia.org/wiki/Sassarese_language}
\item Sample (Italian introduction and letter A) of Bazzoni's Sassarese--Italian phraseological dictionary (2001), listed in the Endangered Languages catalogue: \url{https://www.endangeredlanguages.com/node/105379}
\item Lexilogos page of Sardinian dictionaries, including a Sassarese--Italian dictionary with audio: \url{https://www.lexilogos.com/sarde_dictionnaire.htm}
\item Scholarly philological edition of the Gospel of Matthew in Sassarese (San Matteo), which compares Sassarese with Tempiese: \url{https://www.filologiasarda.eu/files/documenti/pubblicazioni_pdf/cfssanmatteo/02edizione.pdf}
\item Verbix conjugation page: \url{https://docs.verbix.com/Languages/SardinianSassarese}
\end{itemize}

\section{Religious and Community Resources}
The Gospel of Matthew edition above is the only scripture-related item found. No community organisation was identified in the search results.

\section{Digital and Audio Resources}
\begin{itemize}
\item Endangered Languages Project sample page: \url{https://endangeredlanguages.com/lang/5361/samples/9618}
\item The audio dictionary on the Lexilogos page above.
\end{itemize}

\section{Curriculum Development Context}
Sardinian-wide frameworks exist (for example the Limba Sarda Comuna standard), but no Sassarese-specific school materials were found, so any programme would have to build on the Rubattu lessons and the dictionary sources, ideally with local authors. Only a few sources were reached, and more exist in Italian and in print.
\end{document}